However, these positions are continuous but the modules must occupy the $|G|$ integer slots of the super-grid, one module per slot. We therefore snap the spectral embedding onto the grid by solving the classical \emph{assignment problem}: given the chip coordinates $g_1, \ldots, g_{|G|}$ of the super-grid slots, find the bijection $\pi : G \to \{0, \ldots, |G|-1\}$ that minimizes the total squared displacement
\begin{equation}
    \sum_g \bigl\lVert x_g - g_{\pi(g)}\bigr\rVert_2^2.
    \label{eq:assignment}
\end{equation}
Before matching we rescale each axis of the continuous embedding independently so that its range spans the slot bounding box\textemdash necessary because the slot layout is generally not square when $|G|$ is not a perfect square. We then solve~\eqref{eq:assignment} in $O(|G|^3)$ time with the Hungarian algorithm~\cite{hungarian}, which returns the globally optimal assignment through primal-dual iteration on the bipartite graph of modules and slots. Because the eigenvectors $u_2, u_3$ are only unique up to sign and axis labeling, we run the matching across all eight sign flips and axis swaps of $(u_2, u_3)$ and keep the assignment with the lowest cost. The full spectral pipeline\textemdash one eigendecomposition plus one Hungarian matching\textemdash is deterministic, runs in under a second even for $|G|=195$, and delivers the globally optimal continuous L$^2$ layout together with its optimal snap-to-grid rounding.
